% ============================================================================
% Fichier  : partie2/P02_C09_outils_acquisition.tex
% Titre    : Outils et Techniques d'Acquisition
% Auteur   : MINKA MI NGUIDJOÏ Thierry Emmanuel
% Version  : 1.0 -- Juin 2026
% Statut   : RÉDIGÉ
% ============================================================================

\chapter{Outils et Techniques d'Acquisition}
\label{chap:acq}

\epigraph{%
    « By approaching each case methodically, you can evaluate the evidence
    thoroughly and document the chain of evidence. »%
}{--- Amelia Phillips}

\niveauChapitre{Intermédiaire}{%
    Notions de systèmes de fichiers et de ligne de commande recommandées.
    Lecture du Chapitre~\ref{chap:custody} obligatoire~: ce chapitre met
    en \oe uvre les procédures de custody.%
}

\begin{boxobjectifs}
\begin{itemize}
    \item Maîtriser le \termecle{write blocker} matériel et logiciel~:
          principe, activation, vérification.
    \item Conduire une acquisition forensique complète avec FTK~Imager~:
          formats E01, DD~; disque, RAM et mobile.
    \item Comprendre les cinq formats d'image forensique et savoir choisir
          le format adapté au contexte.
    \item Extraire et analyser les données d'un téléphone Android avec ADB.
    \item Conduire une analyse mémoire de base avec Volatility~3.
    \item Analyser une capture réseau avec Wireshark pour identifier
          des communications suspectes.
    \item Appliquer la grille CRO à chaque décision technique d'acquisition.
\end{itemize}
\end{boxobjectifs}

\bigskip

% ============================================================================
\section{Le Write Blocker~: Protection Absolue de l'Evidence}
\label{sec:write-blocker}
% ============================================================================

\subsection{Principe et nécessité}

Le \termecle{write blocker}\index{write blocker} est un dispositif ---
matériel ou logiciel --- qui permet de lire un support numérique en
interdisant toute écriture. Sans write blocker, le simple branchement
d'un disque sur un PC Windows déclenche une cascade invisible de
modifications~: création de fichiers temporaires, mise à jour de la
base de registre, modification des horodatages d'accès NTFS. En
quelques secondes, l'evidence est contaminée.

\begin{equation}
\text{Write blocker} \implies
\forall\, \text{opération d'écriture} \to \texttt{DENIED}
\quad\text{pendant toute la durée de l'acquisition}
\label{eq:write-blocker}
\end{equation}

\subsection{Write blocker matériel}

Le write blocker matériel est un boîtier physique qui s'intercale entre
le support à acquérir et le poste d'acquisition. C'est le standard de
référence~: indépendant du système d'exploitation, il ne peut pas être
contourné par un bug logiciel.

\begin{tableaucours}{p{3.5cm}p{3.5cm}p{5.0cm}}{%
    \tableauHeader{Modèle de référence} &
    \tableauHeader{Interfaces} &
    \tableauHeader{Usage}
}
Tableau Forensics WiebeTech & SATA, IDE, SAS, USB, SD &
    Standard FBI~; \textasciitilde{}~500~USD \\
ICS Image MASSter Solo & SATA, IDE, USB &
    Acquisition autonome sans PC \\
Wiebetech Forensic UltraDock & SATA, USB~3.0, NVMe &
    Portable~; adapté aux SSD modernes \\
\end{tableaucours}
\vspace{-0.8em}
\begin{center}{\small\itshape Tableau~9.1 --- Write blockers matériels de référence}\end{center}

\subsection{Write blocker logiciel}

En l'absence de write blocker matériel --- situation fréquente dans les
unités PNC camerounaises actuelles --- le write blocker logiciel est une
alternative acceptable si correctement documentée.

\begin{boxoutils}{Write blocker logiciel Windows --- activation registre}
\begin{lstlisting}[language=bash_forensique]
@ Activation
reg add HKLM\SYSTEM\CurrentControlSet\Control\StorageDevicePolicies
    /v WriteProtect /t REG_DWORD /d 1 /f

@ Verification
reg query HKLM\SYSTEM\CurrentControlSet\Control\StorageDevicePolicies
@ Attendu : WriteProtect  REG_DWORD  0x1

@ Test validation (doit echouer)
echo test > D:\test_write.txt
@ Attendu : Acces refuse

@ Desactivation apres acquisition
reg add HKLM\SYSTEM\CurrentControlSet\Control\StorageDevicePolicies
    /v WriteProtect /t REG_DWORD /d 0 /f
\end{lstlisting}
\end{boxoutils}

\begin{boxoutils}{Write blocker logiciel Linux}
\begin{lstlisting}[language=bash_forensique]
# Identifier le disque (ne JAMAIS supposer sdb)
lsblk

# Passer en lecture seule AVANT montage
blockdev --setro /dev/sdb
blockdev --getro /dev/sdb    # Attendu : 1

# Montage en lecture seule (double protection)
mount -o ro,noatime /dev/sdb1 /mnt/evidence

# Test validation (doit echouer)
touch /mnt/evidence/test.txt
# Attendu : Read-only file system
\end{lstlisting}
\end{boxoutils}

\begin{boiteRemarque}{Write blocker logiciel vs matériel}
Le write blocker logiciel est acceptable en investigation préliminaire.
En contexte judiciaire (comparution devant le TCS), le write blocker
matériel est recommandé pour éviter toute contestation. Documenter
obligatoirement~: activation + test de vérification (captures d'écran
horodatées dans le formulaire de custody).
\end{boiteRemarque}

% ============================================================================
\section{Formats d'Image Forensique}
\label{sec:formats-image}
% ============================================================================

\begin{tableaucours}{p{1.8cm}p{2.0cm}p{2.2cm}p{5.5cm}}{%
    \tableauHeader{Format} &
    \tableauHeader{Compression} &
    \tableauHeader{Métadonnées} &
    \tableauHeader{Usage recommandé}
}
\textbf{E01} & Oui (LZX) & Riches &
    \textbf{Standard de référence.} Autopsy, FTK, EnCase.
    Préférer pour toute investigation. \\
\textbf{Ex01} & Oui & Riches + AES &
    E01 chiffré~; données sensibles \\
\textbf{DD / Raw} & Non & Aucune &
    Compatible \emph{tous} outils Linux~; précision absolue~;
    taille = taille disque original \\
\textbf{AFF4} & Oui & Très riches + signature &
    Standard émergent~; signé cryptographiquement \\
\textbf{VMDK / VHD} & Variable & Limitées &
    Machines virtuelles~; Autopsy ouvre directement \\
\end{tableaucours}
\vspace{-0.8em}
\begin{center}{\small\itshape Tableau~9.2 --- Formats d'image forensique comparés}\end{center}

\begin{boxdef}
\textbf{Format E01}\index{E01, format}~: format d'Exterro (EnCase). Il
segmente l'image en fichiers configurables (\texttt{.E01}, \texttt{.E02}~\ldots)
et intègre~: hash MD5 et SHA-1 de chaque segment~; métadonnées de l'affaire
(N°~dossier, examinateur, notes)~; vérification d'intégrité automatique
à l'ouverture. C'est le format le plus reconnu devant les tribunaux.
\end{boxdef}

% ============================================================================
\section{FTK Imager~: Acquisition Disque et RAM}
\label{sec:ftk-imager}
% ============================================================================

\subsection{Procédure d'acquisition disque en 7 étapes}

\begin{boxoutils}{FTK Imager --- acquisition disque}
\textbf{Étape~1~:} Activer le write blocker. Vérifier l'espace disponible
sur le support de destination (minimum~: 110\% de la taille source).

\textbf{Étape~2~:} \texttt{File~$\to$ Create Disk Image}. Type de source~:
\begin{itemize}
    \item \textbf{Physical Drive}~: image complète incluant l'espace
          non alloué~; recommandé en investigation forensique.
    \item \textbf{Logical Drive}~: partition uniquement~; plus rapide
          mais peut manquer des artefacts.
\end{itemize}

\textbf{Étape~3~:} Format E01. Options~: segmentation 0 (pas de segmentation)~;
compression~5~; cocher \textit{Verify images after they are created}.

\textbf{Étape~4~:} Métadonnées~:
\begin{lstlisting}[language=bash_forensique]
Case Number    : DOS-2026-034
Evidence Number: EV-001
Unique Description: Dell Inspiron 15 SN-GF8T3Y2 MBONGO
Examiner       : OPJ NKENG Eric - Matricule PNC 4471-B
Notes          : PC allume - RAM capturee prealablement
\end{lstlisting}

\textbf{Étape~5~:} Support de destination \emph{distinct} de la source
(jamais le même disque). Nom sans espaces~:
\texttt{MBONGO\_EV001\_20260315.E01}.

\textbf{Étape~6~:} Lancer. FTK~Imager affiche la progression, les hash
MD5 et SHA-1 en temps réel, et les erreurs de lecture (bad sectors).

\textbf{Étape~7~:} Vérifier que hash source = hash image. Reporter dans
le formulaire de custody. Sceller l'original.
\end{boxoutils}

\begin{boxalert}
\textbf{Erreurs fréquentes~:}
\begin{itemize}
    \item Acquérir sur le même disque~: l'acquisition échoue à mi-chemin.
    \item Oublier les métadonnées~: sans Case Number, pas de valeur probatoire.
    \item Ne pas vérifier les hash~: si source $\neq$ image, recommencer.
    \item Format DD sur 1~To~: 1~To non compressé~; préférer E01
          ($\approx$~50\% compression).
\end{itemize}
\end{boxalert}

\subsection{Acquisition de la mémoire RAM}

La mémoire RAM contient ce qui n'existe nulle part ailleurs~: processus
en cours d'exécution, connexions réseau actives, mots de passe en
mémoire, clés de chiffrement actives. Sa capture requiert que le PC soit
allumé et constitue la \textbf{première action} sur tout système actif
(RFC~3227, rang~2).

\begin{boxoutils}{FTK Imager --- capture mémoire RAM}
\begin{lstlisting}[language=bash_forensique]
# Interface graphique FTK Imager :
# File --> Capture Memory
# Destination : support externe (jamais disque interne)
# Cocher 'Include pagefile' si disponible
# Cliquer Capture Memory

# Resultat : fichier .mem de la taille de la RAM installee
# 16 Go RAM  --> fichier ~16 Go
# 32 Go RAM  --> fichier ~32 Go

# Hash SHA-256 immediatement apres capture :
certutil -hashfile MBONGO_RAM_20260315.mem SHA256
# Reporter dans le formulaire de custody
\end{lstlisting}
\end{boxoutils}

\begin{boiteRemarque}{La RAM modifie la RAM}
La capture de la mémoire modifie légèrement la mémoire~: le processus
FTK~Imager utilise lui-même de la RAM pendant la capture. C'est inévitable
et documenté --- un fait, pas une violation de custody. Documenter dans le
formulaire~: \textit{``Légère modification de la RAM inhérente à l'outil,
conformément à ACPO Principe~2. Capture réalisée avant toute autre action.''}
\end{boiteRemarque}

\subsection{Acquisition Linux avec dd}

Sur Kali Linux ou tout système Linux, l'outil natif \texttt{dd} permet une
acquisition bit-à-bit sans aucun logiciel supplémentaire~:

\begin{lstlisting}[language=bash_forensique,
    caption={Acquisition forensique dd + dcfldd}]
# Identifier le disque (ne jamais supposer sdb)
lsblk -o NAME,SIZE,FSTYPE,MOUNTPOINT

# Hash SHA-256 AVANT acquisition (original)
sha256sum /dev/sdb | tee hash_source_EV001.txt

# Acquisition dd
dd if=/dev/sdb of=/mnt/evidence/EV001.dd
   bs=65536 conv=noerror,sync status=progress
# bs=65536 : blocs 64 Ko (optimal)
# conv=noerror,sync : continue sur erreur, padde zeros

# Alternative dcfldd (hash pendant l'acquisition)
dcfldd if=/dev/sdb of=/mnt/evidence/EV001.dd
   bs=65536 hash=sha256 hashwindow=10M
   hashlog=/mnt/evidence/EV001_hashlog.txt

# Hash SHA-256 APRES (image)
sha256sum /mnt/evidence/EV001.dd | tee hash_image_EV001.txt

# Comparaison (identiques = integrite OK)
diff hash_source_EV001.txt hash_image_EV001.txt
echo "INTEGRITE: $?"    # 0 = OK
\end{lstlisting}

% ============================================================================
\section{Acquisition Mobile~: Android et iOS}
\label{sec:acquisition-mobile}
% ============================================================================

\subsection{Les quatre niveaux d'acquisition mobile}

\begin{tableaucours}{p{2.0cm}p{2.5cm}p{2.5cm}p{5.5cm}}{%
    \tableauHeader{Niveau} &
    \tableauHeader{Technique} &
    \tableauHeader{Données accessibles} &
    \tableauHeader{Conditions}
}
Niveau~1 & Manuelle (écran) & Contacts, SMS, photos visibles &
    Téléphone déverrouillé~; documentation par photo \\
Niveau~2 & Logique (ADB) & Système de fichiers utilisateur &
    Android déverrouillé~; Options développeur activées \\
Niveau~3 & Physique & Mémoire flash complète~; données supprimées &
    Outils spécialisés (Cellebrite, GrayKey) \\
Niveau~4 & Chip-off / JTAG & Extraction directe de la puce &
    Matériel spécialisé~; risque destruction~; dernier recours \\
\end{tableaucours}
\vspace{-0.8em}
\begin{center}{\small\itshape Tableau~9.3 --- Niveaux d'acquisition mobile
(NIST~SP~800-101r1)}\end{center}

\subsection{ADB --- Android Debug Bridge}

\termecle{ADB}\index{ADB} (\textit{Android Debug Bridge}) est l'outil
officiel Google pour communiquer avec un téléphone Android. En forensique,
il constitue le niveau~2~: accès au système de fichiers utilisateur sans
modifier les données.

\textbf{Prérequis~:}
\begin{enumerate}
    \item Mode avion activé (obligatoire --- cf. Ch.~\ref{chap:custody}).
    \item Options développeur~: \texttt{Paramètres~$\to$ À propos~$\to$
          Numéro de build~$\to$ appuyer 7~fois.}
    \item Débogage USB activé~: \texttt{Options développeur~$\to$
          Débogage USB~$\to$ Activé.}
    \item Autorisation sur le téléphone lors du premier branchement.
\end{enumerate}

\begin{boxoutils}{ADB --- commandes forensiques essentielles}
\begin{lstlisting}[language=bash_forensique]
# Verifier la connexion
adb devices
# Attendu : R5CR71XXXX   device

# Identifier l'appareil
adb shell getprop ro.product.model       # Modele
adb shell getprop ro.build.version.release  # Android version
adb shell getprop ro.serialno            # Numero de serie

# Extraction du systeme de fichiers utilisateur
adb pull /sdcard/ ./extraction/sdcard/
# /sdcard : photos, videos, telechargements, backups WhatsApp

# Liste des applications installees
adb shell pm list packages -f > ./extraction/applications.txt

# Extraction SMS (selon version Android)
adb shell content query --uri content://sms/
    > ./extraction/sms.txt

# Contacts
adb shell content query --uri content://contacts/people/
    > ./extraction/contacts.txt

# Journal systeme
adb logcat -d > ./extraction/logcat.txt

# Sauvegarde complete
adb backup -all -apk -shared -f backup_evidence.ab

# Hash SHA-256 de l'extraction
find ./extraction/ -type f -exec sha256sum {} \;
    > hash_extraction.txt
\end{lstlisting}
\end{boxoutils}

\begin{boxCRO}
\textbf{Analyse CRO --- Extraction ADB complète vs ciblée}

L'extraction complète (\texttt{adb pull /sdcard/}) capture tout~:
photos personnelles, messages privés, données de santé.

$\mathcal{C}$~: extraction complète = atteinte maximale à la vie privée~;
extraction ciblée = atteinte proportionnée.

$\mathcal{R}$~: extraction complète préférable pour ne pas manquer
des preuves dans des répertoires inattendus.

$\mathcal{O}$~: le principe de proportionnalité (art.~52
\loicam{2010/012}) exige que la saisie soit limitée au pertinent.

\cro{$\downarrow$ tension à gérer}{$\uparrow$ extraction complète}{$\sim$ justification requise}

\textbf{Règle pratique~:} en contexte judiciaire, documenter pourquoi
une extraction complète était nécessaire (impossibilité de cibler~;
risque de preuves dans des répertoires inattendus).
\end{boxCRO}

\subsection{iOS --- contraintes spécifiques}

\begin{itemize}
    \item Chiffrement hardware depuis iOS~8~; pas d'équivalent ADB~;
    \item Extraction logique via iTunes/Finder uniquement avec code PIN~;
    \item Extraction physique~: outils commerciaux (Cellebrite UFED,
          Magnet AXIOM)~;
    \item Données iCloud~: réquisition à Apple
          (\texttt{legalprocess.apple.com}).
\end{itemize}

En l'absence d'outils spécialisés~: conserver éteint ou mode avion,
documenter l'IMEI, transmettre à un laboratoire accrédité.
\textbf{Ne jamais tenter de brute-forcer le code PIN}~: après
10~tentatives échouées, iOS efface les données.

% ============================================================================
\section{Volatility~3~: Analyse de la Mémoire RAM}
\label{sec:volatility}
% ============================================================================

\subsection{Architecture et principe}

\termecle{Volatility~3}\index{Volatility} est le standard mondial d'analyse
forensique de la mémoire RAM. Il analyse le fichier \texttt{.mem} capturé
avec FTK~Imager et révèle l'état du système \emph{au moment de la capture}~:
processus en cours, connexions réseau, mots de passe, clés de chiffrement.

La version~3 (depuis 2021) utilise une syntaxe unifiée sans profil
obligatoire~: le profil Windows est détecté automatiquement.

\begin{boxoutils}{Volatility~3 --- commandes essentielles}
\begin{lstlisting}[language=bash_forensique]
# -- IDENTIFICATION --
python3 vol.py -f memoire.mem windows.info

# -- PROCESSUS --
# Liste complete (PID, PPID, dates)
python3 vol.py -f memoire.mem windows.pslist

# Arborescence parent-enfant
python3 vol.py -f memoire.mem windows.pstree
# Indicateur suspect : svchost.exe enfant de cmd.exe
# (svchost legitime a toujours services.exe comme parent)

# -- RESEAU --
python3 vol.py -f memoire.mem windows.netscan
# Indicateur suspect : connexion ESTABLISHED vers port 4444

# -- COMMANDES EXECUTEES --
python3 vol.py -f memoire.mem windows.cmdline
# Indicateur suspect : PowerShell encodedcommand base64

# -- FICHIERS EN MEMOIRE --
python3 vol.py -f memoire.mem windows.filescan

# Extraire un fichier depuis la memoire
python3 vol.py -f memoire.mem windows.dumpfiles
    --virtaddr 0xXXXXXXXX

# -- HASHES MOTS DE PASSE --
python3 vol.py -f memoire.mem windows.hashdump

# -- CAPTURES D'ECRAN --
python3 vol.py -f memoire.mem windows.screenshotss
    -D ./screenshots/
\end{lstlisting}
\end{boxoutils}

\subsection{Détecter un processus malveillant~: cas MVOM}

\begin{lstlisting}[language=bash_forensique,
    caption={Analyse Volatility --- Opération MVOM TRANSBOIS SA}]
# Etape 1 : processus
python3 vol.py -f transbois.mem windows.pslist
# Resultat SUSPECT :
# 4892  4891  svchost.exe   2026-03-13 02:17:44
# 4891  4888  cmd.exe       2026-03-13 02:17:42
# ANORMAL : svchost ne doit pas avoir cmd comme parent

# Etape 2 : connexions reseau
python3 vol.py -f transbois.mem windows.netscan | grep 4892
# Resultat :
# ESTABLISHED  10.0.0.15:49212  185.220.XX.XX:4444  4892
# Connexion active vers C2 (port 4444)

# Etape 3 : identifier le fichier malveillant
python3 vol.py -f transbois.mem windows.filescan | grep enc_all
# Resultat :
# \Device\...\AppData\Local\Temp\enc_all.ps1
# Preuve de l'installation au moment de la capture
\end{lstlisting}

% ============================================================================
\section{Wireshark~: Analyse du Trafic Réseau}
\label{sec:wireshark}
% ============================================================================

\subsection{Modes d'utilisation forensique}

\termecle{Wireshark}\index{Wireshark} s'utilise en deux modes~:
\begin{enumerate}
    \item \textbf{Analyse de capture fournie}~: ouvrir un fichier
          \texttt{.pcap} / \texttt{.pcapng} collecté sur la scène.
    \item \textbf{Capture live}~: capturer en temps réel (nécessite
          une position dans le réseau~: SPAN port, mirror, hub).
\end{enumerate}

\subsection{Filtres de référence}

\begin{boxoutils}{Wireshark --- filtres forensiques essentiels}
\begin{lstlisting}[language=bash_forensique]
# Tout le trafic d'une IP
ip.addr == 185.220.XX.XX

# Communications entre deux hotes
ip.addr == 10.0.0.15 && ip.addr == 185.220.XX.XX

# HTTP (non chiffre)
http

# Soumissions de formulaires (mots de passe)
http.request.method == "POST"

# Requetes DNS (domaines consultes)
dns

# Port suspect (Metasploit, shells inverses)
tcp.port == 4444

# Trafic email non chiffre
smtp || pop3 || imap

# Recherche textuelle dans les trames
frame contains "password"
frame contains "Authorization"

# FONCTIONS UTILES :
# Statistics --> Conversations : toutes les paires IP
# Statistics --> Protocol Hierarchy : volumes par protocole
# File --> Export Objects --> HTTP : extraire fichiers telecharges
\end{lstlisting}
\end{boxoutils}

\begin{boxscenario}{Wireshark --- cas WOURI BEC}
La capture réseau SOTRADI.pcap est disponible. L'analyste reconstruit
l'attaque~:
\begin{enumerate}
    \item Filtre \texttt{smtp}~$\to$ email entrant depuis
          \texttt{contact@medequip-portugal.com}.
    \item Clic droit~$\to$ \textit{Follow TCP Stream}~: les en-têtes
          SMTP révèlent l'IP source réelle~: serveur russe.
    \item Filtre \texttt{dns}~$\to$ chercher \texttt{medequip-portugal.com}~:
          la résolution confirme l'enregistrement récent (12/01/2026,
          3~jours avant le paiement).
    \item \textit{Statistics~$\to$ Conversations}~: volumes de transfert
          anormaux vers IPs suspectes.
\end{enumerate}

\textbf{Constatation forensique (sans interprétation)~:} ``Le paquet
\#~14~872 (timestamp 2026-01-15~14h28) contient un email SMTP depuis
\texttt{contact@medequip-portugal.com}, originant depuis l'IP
\texttt{91.XXX.XX.XX} selon le champ \texttt{Received}. Le domaine
\texttt{medequip-portugal.com} a été enregistré le 12/01/2026 selon
les bases WHOIS.''
\end{boxscenario}

% ============================================================================
\section{File Carving~: Récupérer les Fichiers Supprimés}
\label{sec:file-carving}
% ============================================================================

\subsection{Principe des magic bytes}

Le \termecle{file carving}\index{file carving} récupère des fichiers
depuis les données brutes d'un disque en cherchant les
\termecle{magic bytes}\index{magic bytes} --- séquences d'octets
caractéristiques de chaque type de fichier --- sans se fier à la table
de fichiers (qui peut être effacée ou corrompue).

\begin{tableaucours}{p{2.5cm}p{3.5cm}p{5.5cm}}{%
    \tableauHeader{Type} &
    \tableauHeader{Magic bytes (hex)} &
    \tableauHeader{Usage forensique}
}
JPEG & \texttt{FF D8 FF E0} &
    Photos~; contient souvent les métadonnées GPS \\
PDF & \texttt{25 50 44 46} &
    Documents (factures, contrats, bons de commande) \\
ZIP / Office & \texttt{50 4B 03 04} &
    .docx, .xlsx, .pptx (format ZIP interne) \\
PNG & \texttt{89 50 4E 47} &
    Images~; captures d'écran \\
EXE / DLL & \texttt{4D 5A} (``MZ'') &
    Exécutables~; malware supprimé \\
\end{tableaucours}
\vspace{-0.8em}
\begin{center}{\small\itshape Tableau~9.4 --- Magic bytes des types courants}\end{center}

\begin{lstlisting}[language=bash_forensique,
    caption={File carving avec Foremost et Scalpel}]
# Foremost -- types cibles
foremost -t jpg,pdf,docx,xlsx -i evidence.dd -o /output/carving/

# Resultats
ls /output/carving/
# jpg/  pdf/  docx/  audit.txt

cat /output/carving/audit.txt
# nombre de fichiers, taille totale, positions

# Scalpel -- configuration flexible
scalpel -c /etc/scalpel/scalpel.conf evidence.dd
    -o /output/scalpel/

# Rechercher des strings dans l'espace non alloue
strings -a evidence.dd
    | grep -i "momo\|virement\|password\|iban"
    > strings_suspects.txt
\end{lstlisting}

% ============================================================================
\section{Synthèse~: Choisir ses Outils selon la Situation}
\label{sec:choix-outils}
% ============================================================================

\begin{tableaucours}{p{4.5cm}p{3.5cm}p{4.0cm}}{%
    \tableauHeader{Situation} &
    \tableauHeader{Outil} &
    \tableauHeader{Procédure}
}
PC allumé --- capturer la RAM & FTK~Imager &
    \texttt{File~$\to$~Capture Memory} \\
PC éteint --- image disque (Windows) & FTK~Imager &
    \texttt{File~$\to$~Create Disk Image~$\to$~E01} \\
PC éteint --- image disque (Linux) & \texttt{dd} / \texttt{dcfldd} &
    \texttt{dd if=/dev/sdb of=image.dd bs=65536} \\
Android déverrouillé & ADB &
    \texttt{adb pull /sdcard/} + backup \\
iOS & iTunes & Sauvegarde~; réquisition Apple LEA \\
Analyse RAM & Volatility~3 &
    \texttt{windows.pslist}, \texttt{netscan}, \texttt{filescan} \\
Analyse capture réseau & Wireshark &
    Filtres \texttt{http}, \texttt{smtp}, \texttt{tcp.port} \\
Récupérer fichiers supprimés & Foremost / Autopsy &
    \texttt{foremost -t jpg,pdf} ou Deleted Files \\
\end{tableaucours}
\vspace{-0.8em}
\begin{center}{\small\itshape Tableau~9.5 --- Sélection des outils par situation}\end{center}

\separateur

\begin{boxresume}
\textbf{Ce qu'il faut retenir de ce chapitre~:}
\begin{itemize}
    \item \textbf{Write blocker}~: activer AVANT de brancher l'evidence.
          Matériel = gold standard~; logiciel acceptable si documenté
          (activation + test de refus d'écriture).

    \item \textbf{FTK~Imager}~: 7~étapes (write blocker~$\to$ type~$\to$
          format E01~$\to$ métadonnées~$\to$ destination distincte~$\to$
          lancer~$\to$ vérifier hash). RAM~: \texttt{File~$\to$~Capture
          Memory} sur support externe.

    \item \textbf{Format E01}~: standard~; compression~$\approx$~50\%~;
          métadonnées embarquées~; vérification automatique.
          DD~: compatible tous outils mais taille source = taille image.

    \item \textbf{ADB}~: 4 prérequis (mode avion, options développeur,
          débogage USB, autorisation). \texttt{adb pull /sdcard/} +
          backup. Limites~: WhatsApp chiffré sans root. Gérer
          la tension $\mathcal{C}$ vs $\mathcal{R}$ (proportionnalité).

    \item \textbf{Volatility~3}~: trio de base~: \texttt{windows.pslist}
          + \texttt{netscan} + \texttt{filescan}. Indicateurs suspects~:
          \texttt{svchost.exe} enfant de \texttt{cmd.exe}~; connexion
          ESTABLISHED port~4444.

    \item \textbf{Wireshark}~: quatre filtres clés~: \texttt{ip.addr}
          (isoler une IP)~; \texttt{smtp} (emails)~;
          \texttt{http.request.method == "POST"} (formulaires)~;
          \texttt{frame contains "password"}.

    \item \textbf{File carving}~: \texttt{foremost -t jpg,pdf}~;
          magic bytes~; utile quand la table de fichiers est corrompue
          ou délibérément effacée.
\end{itemize}
\end{boxresume}

\bigskip

\begin{boxexo}[title={\ding{45}\quad Exercice~9.1 --- Choix d'outil \niveaubase}]
Pour chacune des situations, indiquez l'outil, la commande exacte et
le format de sortie~:
\begin{enumerate}[(a)]
    \item PC Windows allumé saisi~: capturer la mémoire RAM.
    \item SSD SATA extrait après extinction~: image forensique sous Linux.
    \item Samsung Galaxy A32, Android~11, déverrouillé~: extraire photos
          et liste des applications.
    \item Image DD 500~Go~: récupérer tous les fichiers PDF supprimés.
\end{enumerate}
\end{boxexo}

\begin{boxexo}[title={\ding{45}\quad Exercice~9.2 --- Analyse Volatility \niveaumoyen}]
La sortie de \texttt{python3 vol.py -f incident.mem windows.pstree}
donne (extrait)~:
\begin{lstlisting}[language=bash_forensique]
PID   PPID  Name              CreateTime
1024  788   services.exe
  1284 1024  svchost.exe      (7 instances normales)
4892  4888  svchost.exe       2026-03-13 02:17:44
  4888 4884  cmd.exe          2026-03-13 02:17:42
    4884 4880  powershell.exe 2026-03-13 02:17:40
      4880 4876  explorer.exe 2026-03-13 02:17:38
\end{lstlisting}
\begin{enumerate}[(a)]
    \item Identifiez au moins deux anomalies dans cette arborescence.
    \item Formulez une constatation forensique (sans interprétation)
          pour chaque anomalie.
    \item Quelles commandes Volatility lanceriez-vous ensuite~?
          Dans quel ordre et pourquoi~?
\end{enumerate}
\end{boxexo}

\begin{boxexo}[title={\ding{45}\quad Exercice~9.3 --- Pipeline complet \niveauexpert}]
Vous arrivez sur la scène MVOM~: serveur Dell PowerEdge R740,
Windows~Server~2022, allumé, connecté au réseau. Le ransomware était
actif il y a 20~minutes.
\begin{enumerate}[(a)]
    \item Rédigez le protocole d'acquisition complet dans l'ordre~:
          chaque action, outil, commande exacte, justification forensique.
    \item Le serveur a 128~Go de RAM. Quel support~? Durée estimée~?
          Comment documenter la modification inévitable de la RAM~?
    \item Appliquez la grille CRO à la décision d'effectuer une capture
          Wireshark live de 60~minutes avant d'éteindre le serveur.
    \item Rédigez les deux premières entrées du formulaire de custody.
\end{enumerate}
\end{boxexo}

\bigskip

\begin{boxinfo}
\textbf{Transition.} La Partie~II est maintenant complète~: PIU
(Ch.~\ref{chap:PIU}), Standards (Ch.~\ref{chap:meth-std}),
Custody (Ch.~\ref{chap:custody}) et ce chapitre constituent
le socle opérationnel de tout ce qui suit.

La \textbf{Partie~III} (Chapitres~10--13) approfondit le cadre normatif~:
normes globales ISO et Budapest, législation mondiale comparative,
droit camerounais et africain, déontologie de l'investigateur. Elle
n'est pas un doublon~: elle donne la \emph{légitimité juridique} qui
confère toute leur valeur aux procédures techniques que vous maîtrisez
maintenant.
\end{boxinfo}
